% Einheit 5 von 5 – Dreiecke und Figuren aus Tangente, Normale und Achsen (bank/tangente-normale-schnittwinkel/e5.jsonl, zusammenbau v0.1)
%% TODO Zeitmarke Einheit 5: Marken-Zeile nicht lesbar
%% TODO Zweigzeile Teil 1: was hier gelernt wird (Satz in Schülersprache aus dem Einheitstitel) – Einheit 5
\einheitenkopf[e5]{Einheit 5 von 5 · Dreiecke und Figuren aus Tangente, Normale und Achsen}
\zweigzeile{Abitur GK}
\setcounter{aufgabe}{22}

%% TODO Titel: Ich-kann-Satz fehlt in der Bank (Platzhalter: Kettenname) – Nr. 23
%% TODO Anweisung: ein Satz über der ersten Teilaufgabe fehlt in der Bank – Nr. 23
\begin{aufgabe}{Dreiecke und Figuren}
%% TODO \rechenplatz{n}: Schrittzahl des Grundfalls fehlt in der Bank (nur bei fester Schreibform, 2.3 b)
\begin{teile}
\teil Eine Tangente und die beiden Koordinatenachsen schließen ein Dreieck ein. Gesucht: Wie lang ist sein Rand? Was ist gefragt? \\ \kreuz{Fläche} \\ \kreuz{Umfang} \\ \kreuz{eine Eigenschaft}
\teil Die Gerade $g(x) = -3x + 6$ schließt mit den Koordinatenachsen ein Dreieck ein. Flächeninhalt? $A =$ \leerfeld
\teil Die Tangente an den Graphen von $h(x) = -\frac{4}{x}$ im Punkt $P(-1 | h(-1))$ und die Koordinatenachsen begrenzen ein Dreieck. Flächeninhalt? \hfill (Abitur 2018 LK) \\ $A =$ \leerfeld
\teil $f(x) = (x + 3) \cdot e^{-x}$ mit $f'(x) = -(x + 2) \cdot e^{-x}$. Die Tangente im Punkt $(0 | f(0))$ bildet mit den Koordinatenachsen ein Dreieck. Umfang, und welcher Punkt hat von allen drei Ecken denselben Abstand? \hfill (Abitur 2022 GK) \\ $U \approx$ \leerfeld , Punkt: \leerfeld
\teil $f(x) = \mathrm{ln}(x^2 + 3)$; die Tangente und die Normale an den Graphen im Punkt $B(1 | f(1))$ begrenzen mit der $y$-Achse ein Dreieck. Flächeninhalt? \hfill (Abitur 2017 LK) \\ $A =$ \leerfeld
\teil $f(x) = -\frac{1}{4}x^2 + 4$ – die Tangenten an den Stellen $-2$ und $2$ bilden mit der $x$-Achse ein Dreieck. Umfang? $U \approx$ \leerfeld
\teil $f(x) = 2e^{\frac{1}{2}x} + 1$; die Tangente im Punkt $S(0 | 3)$ begrenzt mit den Koordinatenachsen ein Dreieck. Weise nach, dass es gleichschenklig ist. \hfill (Abitur 2017 LK)
\teil $h(x) = -2x \cdot e^{-0{,}2x}$; $t$ ist die Tangente an den Graphen im Ursprung. $t$, die $x$-Achse und die Gerade $x = b$ mit $b > 0$ schließen eine Fläche von $49$ FE ein. Berechne $b$. \hfill (Abitur 2021 GK) \\ $b =$ \leerfeld
\teil Für $a > 0$ ist $f_a(x) = 2x \cdot (x - a)^2$; $h$ ist die Gerade durch den Ursprung, die senkrecht auf der Tangente an den Graphen im Ursprung steht. Ein Rechteck hat die beiden gemeinsamen Punkte des Graphen mit der $x$-Achse als benachbarte Ecken, eine Diagonale liegt auf $h$. Weise nach, dass sein Flächeninhalt nicht von $a$ abhängt. \hfill (Abitur 2024 LK)
\end{teile}
\end{aufgabe}

%% TODO Titel: Ich-kann-Satz fehlt in der Bank (Platzhalter: Kettenname) – Nr. 24
%% TODO Anweisung: ein Satz über der ersten Teilaufgabe fehlt in der Bank – Nr. 24
\begin{aufgabe}{Dreiecke und Figuren – Fehler finden, Begründen}
\begin{teile}
\teil Ich finde den Fehler: Die Tangente $y = 3x + 6$ bildet mit den Koordinatenachsen ein Dreieck. Tom rechnet: \rechnung{\text{Nullstelle: } x &= -2 \\ A &= \frac{1}{2} \cdot (-2) \cdot 6 = -6}
\teil Begründe, warum das Dreieck aus einer Tangente und den beiden Koordinatenachsen seinen rechten Winkel im Ursprung hat.
\end{teile}
\end{aufgabe}
